
\chapter{Фрод: кто атакует и~как защищаться}\label{ch:fraud}
Универсального сигнала риска не существует, а~атака всё чаще
маскируется под нормальное поведение. Система смотрит на устройство,
историю аккаунта, частоту попыток, маршрут доставки и поведение после
оплаты и~собирает из этих слабых сигналов рабочую оценку риска.
Строгий фильтр повышает цену ошибочного отказа, а~мягкий пропускает
больше фрода. Настройка антифрода~-- настройка этого баланса.

\section{Карта угроз}\label{sec:fraud-threats}

Карточный фрод давно превратился в~разделённый рынок, где одни
участники крадут данные, другие проверяют их на живучесть, третьи
монетизируют уже готовую атаку. Фильтр, хорошо ловящий карточный
перебор, почти ничего не скажет о~злоупотреблении возвратными платежами,
поэтому антифрод начинается с~карты атак.

Самый старый класс атак~-- скимминг. Он остаётся действующим, пока
в~системе где-то жив \textit{fallback} на магнитную полосу. Накладка
на терминале или банкомате считывает данные полосы, украденные дампы
(dumps) уходят на закрытые площадки и превращаются в~клон карты. Чип
EMV резко снизил ценность такого клона, поскольку каждая операция требует
новой криптограммы, которую нельзя просто переписать. Магнитная полоса
исчезла не везде и~не одновременно, и~пока \textit{fallback} жив, скимминг
остаётся действующей угрозой~\cite{emvco-book1}.

Следующий по масштабу класс атак~-- мошенничество без предъявления
карты (card-not-present, CNP) и карточный перебор (carding).
В~электронной коммерции сценарий выглядит знакомо, когда у~атакующего
есть PAN, срок действия и CVV, полученные из утечки, через фишинг,
вредоносное ПО или на теневом рынке. Дальше начинается \textit{carding},
когда серия мелких тестовых списаний проверяет, какие реквизиты ещё живы,
а~уже потом следует крупная покупка~\cite{mc-fraudulent-merchants-2025}.
Атака оставляет плотный повторяющийся рисунок и~хорошо заметна
в~частотных проверках; набор частотных правил срабатывает до того,
как система увидит первую по-настоящему дорогую
операцию~\cite{visa-auth-management-2020}.

Отдельный класс атак~-- захват учётной записи (Account Takeover, ATO).
Злоумышленник входит через уже существующую зону доверия, используя
утёкшие пары логин/пароль (credential stuffing), фишинг,
социальную инженерию, иногда подмену SIM-карты, и наследует всё,
что система прежде считала знакомым~-- историю покупок, сохранённые
карты, адреса доставки. Отдельная транзакция выглядит почти
нормальной, поэтому тревожным сигналом становится резкая смена
устройства, географии, канала входа или адреса
доставки~\cite{mc-fraudulent-merchants-2025}.

Иного рода проблема~-- ложное оспаривание собственной покупки
(friendly fraud), когда никто не крадёт карту и~не ломает аккаунт,
а~держатель сам совершает оплату, получает товар или услугу
и~затем заявляет, что операция была несанкционированной. Это всплывает
в~цифровых товарах, подписках и сервисах с~нематериальной поставкой,
где труднее показать момент потребления. Без связной истории устройства,
аккаунта, использования услуги и~доставки спор превращается в~слово
против слова. Антифрод и~работа с~оспариваниями связаны плотно; следы,
собранные сегодня, понадобятся
в~главе~\ref{ch:disputes}~\cite{mc-first-party-trust}.

Мошенничество \textit{bust-out} строится на терпении, когда сначала
клиент или продавец ведёт себя образцово, набирает доверие, растит
лимиты или оборот, а~затем резко меняет профиль поведения и уводит
деньги из системы. На стороне эмитента это похоже на кредитный обман~--
несколько карт, аккуратная история, постепенный рост лимита и~затем
дефолт. В~эквайринге рисунок другой~-- спокойный старт, быстрый рост
оборота, всё меньше возвратов и внезапное исчезновение продавца.
Коротким окном наблюдения такие случаи не ловятся, потому что их видно
только там, где у~системы есть длинная память о~поведении и полноценный
мониторинг после подключения~\cite{visa-merchant-screening-service,mc-fraudulent-merchants-2025}.

Объект наблюдения меняется вместе с~типом атаки, потому что где-то важнее
устройство, где-то аккаунт, где-то история доставки и потребления,
где-то профиль самого продавца. Абстрактного подозрения недостаточно;
хороший фильтр всегда привязан к~конкретной уязвимости.

\section{Первая линия: пороговые проверки и статические правила}\label{sec:fraud-rules}

Первая линия защиты работает быстрее всех, потому что ей нужно принять
решение до того, как операция успеет развернуться в~ущерб. Эту задачу
решают пороговые проверки (velocity checks) и статические правила. Они
отвечают на простые вопросы~-- не слишком ли много операций идёт
с~одного IP, не появилось ли на одном устройстве подозрительно много
карт, не прыгает ли сумма слишком резко относительно привычной истории
клиента. Правило сравнивает текущую операцию с~заранее заданным
порогом и~за миллисекунды решает, пускать её дальше или нет. Цена
простоты известна заранее~-- жёсткие пороги режут нормальный поток,
а~мягкие пропускают атаку. Правила хороши как первый заслон и непригодны
на роль единственного судьи.

Базовые правила строятся вокруг трёх измерений~-- частота, концентрация
и резкое отклонение от привычки. Сколько операций пришло с~одного IP
за час, сколько попыток сделала одна карта за десять минут, сколько
карт появилось на одном устройстве за сутки, насколько текущая сумма
выбивается из профиля аккаунта. Такие проверки дёшевы, быстро
исполняются и легко объясняются бизнесу, комплаенсу и регулятору.

У~прозрачности есть предел, потому что атакующий умеет дробить поток
под пороги, менять IP, растягивать серию во времени и подмешивать
легитимный шум.
Правило не знает контекста, поэтому крупная покупка в~аэропорту перед
вылетом для одного клиента~-- норма, а~для другого~-- сильная аномалия.
Статические правила работают как грубый фильтр перед более тонкой
моделью оценки.

Особая категория~-- проверки карточного перебора (carding checks).
Злоумышленник не начинает с~дорогой покупки, поэтому сперва идут серии
микросписаний или авторизаций на символические суммы с~одного IP,
устройства или подсети, иногда по десяткам разных карт. Этот
повторяющийся рисунок и~делает \textit{carding} заметным для частотных правил.
Пока атака проверяет, какие реквизиты ещё живы, она сама оставляет
ритм, который полностью скрыть невозможно.

\section{Цифровой слепок устройства и поведенческая аналитика}\label{sec:fraud-fingerprint}

Цифровой слепок устройства (device fingerprinting) -- набор признаков
браузера или телефона, по сочетанию которых система узнаёт конкретное
клиентское окружение. В~него входят заголовок \texttt{User-Agent},
размер экрана, часовой пояс, набор шрифтов, характеристики WebGL
и особенности рендеринга canvas и audio. Ни один признак по отдельности
не решает задачу, но их комбинация даёт устойчивый профиль. Если
браузер скрывает или специально искажает эти атрибуты, точность падает,
но fingerprinting остаётся сильным сигналом
риска~\cite{eckersley-fingerprinting-2010}.

Слепок становится полезнее с~накоплением истории, потому что устройство,
которое месяцами участвует в~нормальных покупках, набирает доверие, тогда
как новый слепок, возникший в~момент дорогой покупки да ещё и~с~новым
адресом доставки, усиливает подозрение. Из таких накопленных связей
строятся программы доказательства участия клиента в~предыдущих
операциях.

Поведенческая аналитика добавляет данные текущей сессии~-- скорость
набора, траекторию мыши, время заполнения формы и~паузы между
действиями. Человек, вводящий свою карту, и~бот, проливающий
украденные реквизиты, оставляют разный темп. Сбор работает пассивно
и~не добавляет трения в~оформление покупки, тогда как браузеры
с~усиленной защитой приватности обедняют сигнал, а~опытный атакующий
имитирует часть поведения. Поведенческие признаки полезны в~связке
со слепком устройства, а~не вместо него.

На мобильных платформах для сопоставимого уровня данных нужен
отдельный SDK, тогда как часть браузеров и расширений специально
искажает атрибуты. Для обычного фродового потока слепок устройства
остаётся полезным, потому что подделать его полностью трудно,
а подделать и при этом не оставить новых аномалий ещё труднее. EMV
3-D Secure выделяет отдельные группы device information, account
information и merchant risk data для риск-решения
ACS~\cite{emvco-3ds-tech-features}.

\section{Машинное обучение в~антифроде: что работает, а~что переоценено}\label{sec:fraud-ml}

Машинное обучение удерживает сразу множество слабых признаков~-- время
суток, сумму, географию, расстояние от предыдущей операции, профиль
продавца, историю устройства. Там, где по отдельности всё выглядит
терпимо, модель замечает редкую комбинацию. При этом модель видит лишь то,
чему её научили, и~стареет вслед за поведением злоумышленников, поэтому
ML работает как тонкая надстройка над правилами, а~не вместо них.

На практике выигрывают самые качественные признаки, и~для
производственного антифрода хватает градиентного бустинга, поскольку он
проще в~обучении, легче объясняется и хорошо работает на табличных данных.
Сами признаки остаются приземлёнными~-- окна по времени, агрегаты по
продавцу, скользящие суммы, расстояние от предыдущей транзакции.
Открытая литература говорит о~том же, что для табличных наборов данных
о~фроде сильными базовыми моделями остаются деревья и градиентный
бустинг, а~качество и~отбор признаков влияют на итоговые \textit{precision}
и \textit{recall} не меньше, чем выбор
архитектуры~\cite{fraud-feature-selection-2024,fraud-lightgbm-2024}.

Конкретный вектор признаков для одной CNP-транзакции выглядит
примерно так. Сумма транзакции; отношение суммы к~среднему чеку
держателя за 90~дней; количество транзакций этой карты за последний
час и~за сутки; расстояние в~километрах между текущим IP
и IP предыдущей операции; время с~предыдущей операции
в~минутах; возраст аккаунта у~продавца в~днях; совпадение
страны BIN и~страны IP; hash устройства совпадает с~ранее
виденным; MCC ТСП; результат CVV2-проверки; результат
AVS; количество отказов по этой карте за 24~часа; количество уникальных
продавцов за последний час; время суток, нормализованное
к~часовому поясу держателя. Набор из 15--25 таких признаков
покрывает большинство известных атакующих паттернов.

Модель возвращает скор~-- условно от 0 до 1\,000, а~пороги задаются
правилами после скоринга. При низких значениях транзакция проходит
без дополнительной проверки, средние уходят в 3DS-challenge или
ручную очередь, высокие~-- в~автоматический отказ. Пороги калибруются
под целевой false positive rate, а конкретные числа подбираются
эмпирически по собственным данным оператора. Универсальных значений,
применимых к~произвольному потоку, не существует.

Главная производственная проблема~-- дрейф концепции (concept drift).
Мошенники подстраиваются под то, что модель уже научилась ловить,
и~через несколько недель вчерашний сигнал перестаёт работать, а~модель,
давно не переобучавшаяся, деградирует. Хорошо построенная многослойная
защита от фрода устроена как непрерывный процесс с~контролем качества
на скользящем окне, регулярным переобучением и сравнением новой и~старой
версии на отложенных данных~\cite{concept-drift-survey-2014}.

Ещё одна ловушка~-- дисбаланс классов: в~реальном потоке легитимных
операций гораздо больше, чем мошеннических, поэтому простая точность
ничего не говорит о~качестве. Модель, всегда отвечающая
\enquote{легитимно}, формально выглядит почти идеальной и фактически
бесполезна. В~антифроде смотрят на precision, recall и~их компромисс
в~метриках F1 и PR-AUC и оценивают, какую цену система платит за
пропущенный фрод и~за ошибочные
блокировки~\cite{roc-pr-imbalanced-2024}.

Объяснимость в~антифроде не факультативна, потому что клиентскому сервису,
ручным аналитикам и бизнесу нужна понятная причина отказа, а~ответ
в~духе \enquote{модель сказала нет} не помогает ни восстановить
платёж, ни улучшить систему. В~производстве используют SHAP и LIME,
показывающие вклад отдельных признаков в~итоговое
решение~\cite{lime-2016,shap-2017}.

\section{Гибридные системы: правила и~модель}\label{sec:fraud-hybrid}

В~проде побеждает гибридная схема, в~которой правила быстро отсекают
грубые атаки, модель аккуратнее работает с~пограничными случаями,
а завершающие правила переводят численную оценку в~конкретное
действие. Ни один из трёх этапов по отдельности задачу не закрывает; вместе они
дают скорость, чувствительность и управляемость.

Такая защита состоит из трёх этапов, и~сначала идёт предварительный
отсев (pre-screening), при котором быстрые правила блокируют
очевидные случаи~-- карточный перебор, карту из стоп-листа, явное
превышение порога частоты. Затем включается скоринговый модуль
(scoring engine), и~модель выдаёт оценку риска для того, что прошло
первичный фильтр. После этого вступают правила после скоринга
(post-scoring rules), где высокие значения ведут к~блокировке,
средние требуют дополнительной верификации или проверки через 3DS,
а~низкие проходят дальше. Важны не названия этапов, а распределение
ответственности~-- первый этап ловит очевидное, второй сравнивает
сложные паттерны, третий принимает операционное решение.

Требования по латентности для каждого этапа различаются. Правила
pre-screening отрабатывают за единицы миллисекунд, потому что любая
задержка напрямую увеличивает время ответа на авторизацию. ML-скоринг
допускает 20--50 мс, поскольку включается только после быстрого
отсева. Оркестратор решений, переводящий скор в~действие, не требует
вычислений и~задержки практически не добавляет, ограничиваясь
порогами и~набором post-scoring правил.

Правила остаются в~системе даже после того, как модель показывает
хорошее качество, по двум причинам. Во-первых, регуляторные требования
напрямую обязывают блокировать определённые категории операций по
формализованному признаку, и~модель явного правила не заменяет.
Во-вторых, при разборе возвратного платежа или ответе на запрос
регулятора нужна объяснимая причина отказа~-- \enquote{модель так
решила} не подходит, а \enquote{транзакция превысила порог по~частоте}
документируется и проверяется.

\begin{figure}[tbp]
  \centering
  \includefiguresvg[width=.45\textwidth]{assets/figures/ch20-fraud/fraud-hybrid}
  \caption{Гибридная система антифрода: правила и~скоринг.}\label{fig:fraud-hybrid}
\end{figure}

Замкнутая обратная связь обязательна. Исходы транзакций, ручные
проверки, оспаривания и подтверждённые случаи фрода возвращаются
и в~обучающую выборку, и~в~набор правил. Иначе система слепнет,
поскольку вчерашние атаки она ещё узнаёт, а сегодняшние выглядят
новым нормальным поведением. Этот эффект~-- concept drift~-- требует
регулярного переобучения модели на свежих данных. Гибридная схема
работает, пока умеет учиться на собственных ошибках.

Выбор архитектуры антифрода зависит от объёма и~зрелости, а~не от
единого порога транзакций в~месяц. На малом потоке собственная
ML-модель нецелесообразна, потому что обучающая выборка мала, concept
drift незаметен, а стоимость содержания модели сопоставима с~потерями
от фрода~-- достаточно набора правил, предоставляемых PSP. На среднем
потоке имеет смысл гибридная схема, в~которой готовая модель от PSP
или вендора дополняется кастомными правилами под специфику бизнеса.
На больших объёмах собственная модель окупается, поскольку данных
достаточно для обучения, concept drift заметен в~реальном времени,
а~контроль над порогами даёт прямой финансовый эффект. Конкретные
числовые пороги перехода между этими режимами оператор устанавливает
по собственным метрикам потерь от фрода и стоимости владения моделью.

\section{Скоринг и 3DS RBA}\label{sec:fraud-3ds-rba}

Риск-ориентированная аутентификация (Risk-Based Authentication, RBA)
в 3DS v2 (глава~\ref{ch:3ds}) определяет, требовать ли дополнительную
проверку или разрешить авторизацию без подтверждения. Задача
продавца~-- передать эмитенту контекст, который сделает решение
точнее, а~не шумнее.

В 3DS v2 запрос аутентификации несёт набор структурированных данных
от продавца и 3DS Server.
EMVCo отдельно выделяет cardholder account information,
merchant risk indicator, shipping and delivery information
и device information, из которых ACS собирает контекст вместе
со своей собственной моделью~\cite{emvco-3ds,emvco-3ds-tech-features}.

Связка устроена так, что антифрод продавца вычисляет оценку риска
и~передаёт её в 3DS-запросе через релевантные атрибуты, причём при
высоком риске просит более строгую проверку, а~при низком даёт
эмитенту основания пройти без дополнительного подтверждения. ACS
сопоставляет сигнал со своей моделью и принимает собственное решение.
Это кооперативный протокол, в~котором ни продавец, ни эмитент не
управляют исходом единолично, но качественный контекст заметно
повышает вероятность правильного выбора.

На стороне сетей и~сетевых провайдеров давно работают собственные
сервисы риск-оценки~-- у Visa это VAA/VCAS, у Mastercard Decision
Intelligence. Они добавляют эмитенту ещё один уровень сигнала,
помогающий не путать легитимный поток с~серией тестовых и атакующих
попыток. Для продавца это ещё один вход в~общей авторизационной
логике~\cite{visa-auth-management-2020,mc-decision-intelligence}.

\section{Ложноположительное срабатывание как скрытая потеря}\label{sec:fraud-false-positive}

Идеального фильтра не бывает, поскольку система, не пропускающая ни
одной атаки, режет и~слишком много нормальных операций. Настройка антифрода~--
выбор между двумя потерями, ложноположительным срабатыванием (false
positive) и ложноотрицательным срабатыванием (false negative).

Большинство команд эмоционально чувствительны прежде всего
к~пропущенному фроду, потому что каждая мошенническая операция~--
прямой убыток, часто заканчивающийся возвратным платежом. У~ошибочной
блокировки цена тоже есть, и~она прячется в~операционных расходах.

По данным ЕЦБ, общий уровень карточного
фрода в SEPA в 2021 году составил около 0,028\,\% от объёма
карточных операций, причём около 84\,\% этих потерь пришлось
именно на CNP-канал~\cite{ecb-card-fraud-2021}. Для
российского рынка сопоставимые агрегированные показатели
публикует ФинЦЕРТ Банка России в \enquote{Обзоре операций,
совершённых без добровольного согласия клиентов финансовых
организаций}, причём документ охватывает все ЭСП, не только
карты~\cite{cbr-fincert-fraud-overview}.
Мониторинговые программы Visa и Mastercard начинают обращать
внимание на ТСП при уровне фрода выше определённых порогов
(точные значения сверяются по действующим редакциям VAMP/ECP).
Универсальных публичных бенчмарков по false positive rate
для ML-скоринга и rule-only систем нет; разница между ними
измеряется на собственном потоке оператора. Каждый лишний
процент ложноположительных срабатываний -- это примерно
$N \times \mathrm{AOV}$ потерянной выручки в~день, где
$N$ -- объём транзакций и $\mathrm{AOV}$ -- средний чек.

Когда легитимная транзакция ошибочно блокируется, клиент обращается
в~поддержку, а~продавец теряет текущую оплату и~будущие покупки этого
человека. Поддержка обходится дороже, если отказ произошёл при первой
покупке, поскольку такой клиент не возвращается. Ложноположительные
срабатывания считают в~абсолютных
деньгах~\cite{visa-auth-management-2020}.

ROC- и PR-кривые помогают увидеть компромисс на графике. В детекции
фрода точность (precision) и~полнота (recall) считаются для
класса мошеннических операций, а~не для всего потока. Двигаясь по
кривой, команда меняет порог решения, и~чем ниже порог, тем больше
блокировок и меньше мошенничества, но больше ошибочных отказов;
выше порог~-- обратный эффект. Оптимальная точка зависит от экономики
продукта. Для цифровых товаров с~низкой себестоимостью доставки фильтр
строже, а~для физической доставки и~дорогого возврата порог приходится
смягчать. Оптимальный порог~-- экономическое решение, а~не абстрактный
максимум AUC.

\practicenote{%
  Стоимость false positive считают явно, складывая упущенную транзакцию,
  риск потерять клиента и~расходы поддержки. Такой расчёт показывает,
  что система настроена агрессивнее, чем бизнесу выгодно.

}
